\documentclass[12pt,a4paper]{article}
\usepackage[utf8]{inputenc}
\usepackage[T1]{fontenc}
\usepackage[spanish]{babel}
\usepackage{geometry}
\usepackage{xcolor}
\usepackage{listings}
\usepackage{booktabs}
\usepackage{hyperref}
\usepackage{fancyhdr}
\usepackage{titlesec}
\usepackage{array}
\usepackage{longtable}
\usepackage{parskip}
\usepackage{mdframed}
\usepackage{enumitem}
\usepackage{multicol}

\geometry{top=2.5cm, bottom=2.5cm, left=3cm, right=3cm}

\definecolor{codeblue}{RGB}{41, 98, 255}
\definecolor{codegray}{RGB}{128, 128, 128}
\definecolor{codegreen}{RGB}{38, 166, 65}
\definecolor{bgcode}{RGB}{245, 247, 250}
\definecolor{accent}{RGB}{0, 112, 243}
\definecolor{headerblue}{RGB}{15, 30, 70}
\definecolor{taggray}{RGB}{90, 90, 110}

\lstset{
  backgroundcolor=\color{bgcode},
  basicstyle=\ttfamily\footnotesize,
  breaklines=true,
  captionpos=b,
  commentstyle=\color{codegray},
  keywordstyle=\color{codeblue}\bfseries,
  stringstyle=\color{codegreen},
  numberstyle=\tiny\color{codegray},
  numbers=left,
  numbersep=8pt,
  frame=single,
  framesep=4pt,
  rulecolor=\color{codegray!40},
  tabsize=2,
  showstringspaces=false,
  xleftmargin=12pt,
}

\titleformat{\section}{\Large\bfseries\color{headerblue}}{}{0em}{\thesection\quad}[\titlerule]
\titleformat{\subsection}{\large\bfseries\color{accent}}{}{0em}{\thesubsection\quad}
\titleformat{\subsubsection}{\normalsize\bfseries}{}{0em}{\thesubsubsection\quad}

\pagestyle{fancy}
\fancyhf{}
\fancyhead[L]{\small\textcolor{codegray}{SaukCode.cl --- Documentación Técnica}}
\fancyhead[R]{\small\textcolor{codegray}{Vol. II: Componentes}}
\fancyfoot[C]{\small\textcolor{codegray}{\thepage}}
\renewcommand{\headrulewidth}{0.4pt}

\mdfdefinestyle{infobox}{
  linecolor=accent, linewidth=1.5pt, backgroundcolor=accent!5,
  innertopmargin=8pt, innerbottommargin=8pt,
  innerleftmargin=12pt, innerrightmargin=12pt, roundcorner=4pt,
}

\mdfdefinestyle{propbox}{
  linecolor=codegreen!60, linewidth=1pt, backgroundcolor=codegreen!5,
  innertopmargin=6pt, innerbottommargin=6pt,
  innerleftmargin=10pt, innerrightmargin=10pt, roundcorner=3pt,
}

\newcommand{\prop}[3]{%
  \texttt{#1} \textit{(#2)} \textemdash\ #3\\%
}

\hypersetup{colorlinks=true, linkcolor=accent, urlcolor=codeblue}

\begin{document}

% Portada
\thispagestyle{empty}
\begin{center}
  \vspace*{3cm}
  {\Huge\bfseries\color{headerblue} SaukCode.cl}\\[0.5cm]
  {\Large\color{accent} Documentación Técnica}\\[0.3cm]
  {\large\color{codegray} Volumen II: Catálogo de Componentes}\\[2cm]
  \rule{10cm}{1.5pt}\\[0.5cm]
  {\normalsize Layout \textbullet\ MDX \textbullet\ Features \textbullet\ UI \textbullet\ Providers}\\[0.5cm]
  \rule{10cm}{1.5pt}\\[2cm]
  {\small\color{codegray} Generado: Mayo 2026}
\end{center}
\newpage

\tableofcontents
\newpage

% =====================================================
\section{Visión General de Componentes}
% =====================================================

Los componentes están organizados en 6 categorías bajo \texttt{src/components/}:

\begin{table}[h]
\centering
\begin{tabular}{lll}
\toprule
\textbf{Directorio} & \textbf{Tipo} & \textbf{Cantidad} \\
\midrule
\texttt{layout/} & Estructura global de página & 4 \\
\texttt{common/} & Reutilizables genéricos & 6+ \\
\texttt{ui/} & Elementos de interfaz & 3 \\
\texttt{mdx/} & Específicos de contenido MDX & 15+ \\
\texttt{features/} & Por sección del sitio & 20+ \\
\texttt{providers/} & Proveedores de contexto & 2 \\
\bottomrule
\end{tabular}
\caption{Categorías de componentes}
\end{table}

% =====================================================
\section{Componentes de Layout}
% =====================================================

Estos componentes forman el esqueleto visual de todas las páginas.

\subsection{Navbar.jsx}

Barra de navegación principal con soporte para menú móvil y toggle de tema.

\begin{mdframed}[style=propbox]
\textbf{Características:}\\
\prop{darkMode toggle}{boolean state}{Alterna clases \texttt{dark-mode-override}/\texttt{light-mode-override} en \texttt{<html>} y persiste en \texttt{localStorage}}
\prop{menuOpen}{boolean state}{Controla el panel lateral deslizante en móvil}
\prop{isScrolled}{boolean state}{Activa sombra/blur cuando \texttt{scrollY > 10px}}
\end{mdframed}

\textbf{Animaciones}: Usa \texttt{menuSlide} (spring desde derecha) y \texttt{menuItemVariants} (stagger) de \texttt{utils/animations.js}.

\textbf{Comportamiento de tema}:
\begin{lstlisting}[language=JavaScript]
const toggleTheme = () => {
  const isDark = document.documentElement.classList
    .contains('dark-mode-override');
  if (isDark) {
    document.documentElement.classList.remove('dark-mode-override');
    document.documentElement.classList.add('light-mode-override');
    localStorage.setItem('theme', 'light');
  } else {
    document.documentElement.classList.remove(
      'dark-mode', 'light-mode-override'
    );
    document.documentElement.classList.add('dark-mode-override');
    localStorage.setItem('theme', 'dark');
  }
};
\end{lstlisting}

\subsection{Breadcrumb.jsx}

Navegación de migas de pan dinámica basada en la URL actual.

\begin{mdframed}[style=propbox]
\textbf{Características:}\\
\prop{usePathname()}{Next.js hook}{Lee la ruta actual del App Router}
Genera automáticamente las migas dividiendo la URL por \texttt{/} y formateando cada segmento (capitalización, reemplazo de guiones)
\end{mdframed}

\subsection{Footer.jsx}

Pie de página con información del sitio y enlaces.

\subsection{PageTransition.jsx}

Wrapper de transición de página. Envuelve el contenido en un elemento \texttt{motion.div} con variante \texttt{pageTransition} de Framer Motion.

\begin{lstlisting}[language=JavaScript]
// src/app/template.js — envuelve todas las páginas
export default function Template({ children }) {
  return <PageTransition>{children}</PageTransition>;
}
\end{lstlisting}

% =====================================================
\section{Componentes MDX}
% =====================================================

Estos componentes están disponibles dentro de cualquier archivo \texttt{.mdx} del proyecto. Se mapean en \texttt{src/mdx-components.jsx} y en el \texttt{CustomMDXProvider}.

\subsection{Componentes de Texto (\texttt{mdx/text/})}

\subsubsection{Headings (H1--H6)}

Componentes personalizados para todos los niveles de encabezado. Añaden estilos consistentes con soporte de tema oscuro/claro.

\subsubsection{Paragraph}

Párrafo con tipografía mejorada. Aplica \texttt{line-height} y espaciado optimizados para lectura.

\subsubsection{Lists (Ol, Ul)}

Listas ordenadas y no ordenadas con estilos personalizados (marcadores, indentación, espaciado).

\subsection{Componentes de Layout (\texttt{mdx/layout/})}

\subsubsection{SmartTable.jsx}

Tabla avanzada con soporte de fórmulas matemáticas en celdas.

\begin{mdframed}[style=propbox]
\textbf{Props:}\\
\prop{columns}{Array\textless\{key, label, formula?\}\textgreater}{Definición de columnas; si \texttt{formula=true} el contenido se renderiza con KaTeX}
\prop{data}{Array\textless Object\textgreater}{Filas de datos}
\prop{className}{string?}{Clase CSS adicional}
\end{mdframed}

Ejemplo de uso en MDX:
\begin{lstlisting}[language=JavaScript]
<SmartTable
  columns={[
    { key: 'formula', label: 'Fórmula', formula: true },
    { key: 'desc', label: 'Descripción' }
  ]}
  data={[
    { formula: '\\frac{n!}{k!(n-k)!}', desc: 'Combinaciones' }
  ]}
/>
\end{lstlisting}

\subsubsection{Blockquote.jsx}

Cita visual con borde lateral de color y fondo diferenciado. Soporta anidamiento.

\subsubsection{Solution.jsx}

Contenedor colapsable para mostrar soluciones. Usa el elemento nativo \texttt{<details>} envuelto en estilos personalizados.

\begin{lstlisting}[language=HTML]
<Solution>
  El resultado es **42**. El algoritmo tiene complejidad O(n log n).
</Solution>
\end{lstlisting}

\subsubsection{Quiz.jsx y QuizQuestion.jsx}

Sistema de preguntas de opción múltiple interactivo.

\begin{mdframed}[style=propbox]
\textbf{QuizQuestion Props:}\\
\prop{question}{string}{Texto de la pregunta}
\prop{options}{Array\textless string\textgreater}{Lista de opciones de respuesta}
\prop{answer}{number}{Índice (0-based) de la respuesta correcta}
\prop{explanation}{string?}{Explicación que aparece al responder}
\end{mdframed}

\begin{lstlisting}[language=HTML]
<Quiz>
  <QuizQuestion
    question="¿Cuál es la complejidad de quicksort en promedio?"
    options={['O(n)', 'O(n log n)', 'O(n²)', 'O(log n)']}
    answer={1}
    explanation="Quicksort tiene complejidad promedio O(n log n)"
  />
</Quiz>
\end{lstlisting}

\subsubsection{MultiColumn.jsx}

Divide el contenido en múltiples columnas CSS.

\begin{mdframed}[style=propbox]
\prop{columns}{2 | 3}{Número de columnas}
\prop{children}{ReactNode}{Contenido a distribuir en columnas}
\end{mdframed}

\subsection{Componentes de Código (\texttt{mdx/code/})}

\subsubsection{CodeBlock.jsx}

Bloque de código con resaltado de sintaxis, número de líneas y botón de copia.

\begin{mdframed}[style=propbox]
\textbf{Props:}\\
\prop{language}{string}{Lenguaje de programación (python, javascript, etc.)}
\prop{filename}{string?}{Nombre del archivo a mostrar en header}
\prop{showLineNumbers}{boolean?}{Muestra números de línea (default: true)}
\prop{children}{string}{Código fuente}
\end{mdframed}

\subsubsection{CodeComponents.jsx}

Mapeo de elementos de código inline (\texttt{<code>}) a estilos personalizados. Detecta si el código es inline o dentro de un bloque y aplica estilos distintos.

\subsubsection{Terminal.jsx}

Simula una terminal de línea de comandos con prompt visual.

\subsection{Componentes de Media (\texttt{mdx/media/})}

\subsubsection{SmartLink.jsx}

Link inteligente que detecta si la URL es interna o externa:
\begin{itemize}
  \item \textbf{Interna} (\texttt{/}...): Usa \texttt{<Link>} de Next.js para navegación SPA
  \item \textbf{Externa} (http/https): Agrega \texttt{target="\_blank"} y \texttt{rel="noopener noreferrer"}
\end{itemize}

\begin{lstlisting}[language=HTML]
<SmartLink href="/inacap/ti3v53/lectura10">Ver lectura 10</SmartLink>
<SmartLink href="https://example.com">Recurso externo</SmartLink>
\end{lstlisting}

\subsubsection{SmartFigure.jsx}

Wrapper de imágenes con funcionalidades avanzadas.

\begin{mdframed}[style=propbox]
\textbf{Props:}\\
\prop{lightbox}{boolean?}{Activa vista ampliada al hacer clic}
\prop{centered}{boolean?}{Centra la imagen horizontalmente}
\prop{width}{string?}{Ancho máximo de la imagen}
\prop{children}{ReactNode}{Imagen + caption opcional}
\end{mdframed}

\begin{lstlisting}[language=HTML]
<SmartFigure lightbox={true} centered={true}>
  <img src="/diagrama.png" alt="Diagrama de flujo" />
  Figura 1: Diagrama de flujo del algoritmo
</SmartFigure>
\end{lstlisting}

\subsubsection{Iframe.jsx}

Wrapper seguro para iframes (videos, demos). Añade responsive wrapper con aspect ratio correcto.

\subsection{Utilidades MDX (\texttt{mdx/utils/})}

\subsubsection{slugify.js}

Convierte texto a slug URL-friendly para los IDs de encabezados y el Table of Contents.

% =====================================================
\section{Componentes de Features (por Sección)}
% =====================================================

\subsection{INACAP (\texttt{features/inacap/})}

\subsubsection{SubjectsGrid.jsx}

Grid de tarjetas de asignaturas, organizado por semestre. Muestra el código, nombre, semestre y estado de completitud de cada asignatura.

\subsubsection{SemesterTabs.jsx}

Pestañas de navegación para filtrar asignaturas por semestre. Estado interno de pestaña activa.

\subsubsection{Calendar.jsx}

Calendario académico visual con eventos de clases, evaluaciones y feriados.

\begin{mdframed}[style=propbox]
\textbf{Props:}\\
\prop{events}{Array\textless\{date, title, type\}\textgreater}{Lista de eventos del calendario}
\prop{month}{number}{Mes a mostrar (1-12)}
\prop{year}{number}{Año a mostrar}
\end{mdframed}

\subsubsection{ProfessorCard.jsx}

Tarjeta con información del docente: nombre, especialidad, contacto.

\subsubsection{InacapStats.jsx}

Estadísticas de progreso de la asignatura: clases completadas, evaluaciones, porcentaje de avance.

\subsection{Bootcamp (\texttt{features/bootcamp/})}

\subsubsection{CftTable.jsx}

Tabla expandible de módulos para cursos CódigoFacilito. Soporta:
\begin{itemize}
  \item Estructura de módulos con clases anidadas
  \item Links a clases y ejercicios
  \item Estado de completitud por clase
  \item Paginación de módulos
\end{itemize}

\subsubsection{UddDetail.jsx}

Vista de detalle para cursos UDD con diseño específico de esa plataforma.

\subsection{E-Learning (\texttt{features/elearning/})}

\subsubsection{CourseTable.jsx}

Tabla de cursos para una plataforma específica. Muestra título, plataforma, estado y enlace.

\subsubsection{SectionTable.jsx}

Vista de secciones dentro de un curso e-learning.

\subsubsection{CompactPlatformGrid.jsx}

Grid compacto para seleccionar entre las 5 plataformas disponibles.

\subsection{Idiomas (\texttt{features/idioma/})}

\subsubsection{LanguagePage.jsx}

Página principal de un idioma con estadísticas generales y acceso a módulos.

\subsubsection{ModuleTable.jsx}

Tabla de módulos y lecciones para cursos de idiomas.

\subsection{LeetCode (\texttt{features/leetcode/})}

\subsubsection{ProblemTable.jsx}

Tabla interactiva de problemas LeetCode. Integra los tres hooks de tabla.

\begin{mdframed}[style=propbox]
\textbf{Funcionalidades:}\\
--- Búsqueda por título, ID o tópico con normalización de texto (tildes)\\
--- Filtro por dificultad (Fácil / Medio / Difícil)\\
--- Filtro por tópico (array, string, graph, etc.)\\
--- Ordenamiento por ID, título, tópico o dificultad\\
--- Paginación con 20 problemas por página
\end{mdframed}

\subsubsection{ProblemStats.jsx}

Estadísticas globales: total de problemas, distribución por dificultad, porcentajes completados.

\subsection{CodeVault (\texttt{features/codevault/})}

\subsubsection{CodeTable.jsx}

Tabla de algoritmos y estructuras de datos con filtros por categoría y búsqueda.

\subsection{Home (\texttt{features/home/})}

\subsubsection{SectionCards.jsx}

Cards de navegación a las secciones principales del sitio en la página de inicio.

% =====================================================
\section{Componentes de UI}
% =====================================================

\subsection{Card.jsx}

Tarjeta base reutilizable con:
\begin{itemize}
  \item Efecto hover (\texttt{hoverLift} de Framer Motion)
  \item Soporte de tema oscuro/claro vía CSS modules
  \item Composición flexible (cualquier children)
\end{itemize}

\subsection{CardCarousel.jsx}

Carrusel de tarjetas para la página principal. Usa \texttt{staggerContainer} con \texttt{staggerItem} para animar la entrada de cada tarjeta.

% =====================================================
\section{Componentes Comunes}
% =====================================================

\subsection{MDXWrapper.jsx}

Wrapper que aplica estilos base al contenido MDX. Establece ancho máximo, tipografía y espaciado.

\subsection{TableOfContents.jsx}

Tabla de contenidos flotante generada automáticamente a partir de los headings del MDX.

\begin{mdframed}[style=propbox]
Escanea el DOM del documento en busca de elementos \texttt{h2} y \texttt{h3} para construir el índice. Usa \texttt{IntersectionObserver} para destacar la sección visible.
\end{mdframed}

\subsection{ScrollReveal.jsx}

Wrapper de animación de scroll. Aplica una variante de Framer Motion cuando el elemento entra al viewport.

\begin{lstlisting}[language=JavaScript]
<ScrollReveal variant="fadeInUp">
  <MiComponente />
</ScrollReveal>
\end{lstlisting}

\subsection{Componentes de Tabla (\texttt{common/table/})}

\subsubsection{PaginationControls.jsx}

Controles de paginación numerados.

\begin{mdframed}[style=propbox]
\textbf{Props:}\\
\prop{currentPage}{number}{Página actual}
\prop{totalPages}{number}{Total de páginas}
\prop{onPageChange}{function}{Callback al cambiar página}
\prop{getPageNumbers}{function}{Retorna array de números de página visibles}
\prop{maxVisible}{number?}{Máximo de botones de página (default: 5)}
\end{mdframed}

\subsubsection{SortableHeader.jsx}

Encabezado de columna de tabla con indicadores de ordenamiento (↑ ↓ ↕).

\begin{mdframed}[style=propbox]
\textbf{Props:}\\
\prop{columnKey}{string}{Clave de la columna}
\prop{label}{string}{Texto a mostrar}
\prop{sortConfig}{\{key, direction\}}{Estado actual de ordenamiento}
\prop{onSort}{function}{Callback al hacer clic}
\end{mdframed}

\subsubsection{FilterControls.jsx}

Panel de filtros con campos de búsqueda y selects.

\begin{mdframed}[style=propbox]
\textbf{Props:}\\
\prop{searchTerm}{string}{Término de búsqueda actual}
\prop{onSearchChange}{function}{Actualiza búsqueda}
\prop{topicFilter}{string}{Filtro de tópico activo}
\prop{onTopicChange}{function}{Actualiza filtro de tópico}
\prop{difficultyFilter}{string}{Filtro de dificultad activo}
\prop{onDifficultyChange}{function}{Actualiza dificultad}
\prop{allTopics}{Array\textless string\textgreater}{Opciones para el select de tópicos}
\prop{onClear}{function}{Limpia todos los filtros}
\end{mdframed}

% =====================================================
\section{Providers}
% =====================================================

\subsection{ThemeProvider.jsx}

Proveedor ligero que expone el estado del tema si se necesita desde componentes hijo (no usado globalmente — la lógica principal vive en \texttt{Navbar.jsx}).

\subsection{CustomMDXProvider.jsx}

Proveedor de componentes MDX con soporte de tema dinámico.

\begin{lstlisting}[language=JavaScript, caption={CustomMDXProvider.jsx completo}]
'use client';
import { MDXProvider } from '@mdx-js/react';

export function CustomMDXProvider({ children }) {
  const [isDarkMode, setIsDarkMode] = useState(false);

  useEffect(() => {
    const checkDarkMode = () => {
      const isDark =
        document.documentElement.classList.contains('dark-mode') ||
        document.documentElement.classList.contains('dark-mode-override');
      setIsDarkMode(isDark);
    };
    checkDarkMode();
    const observer = new MutationObserver(checkDarkMode);
    observer.observe(document.documentElement, {
      attributes: true, attributeFilter: ['class']
    });
    return () => observer.disconnect();
  }, []);

  const components = {
    details: props => <details {...props} />,
    summary: props => (
      <summary style={{ color: isDarkMode ? '#ff6b6b' : 'red', ... }} />
    )
  };

  return <MDXProvider components={components}>{children}</MDXProvider>;
}
\end{lstlisting}

El \texttt{MutationObserver} garantiza que el componente reaccione en tiempo real a cambios de tema sin necesidad de un sistema de estado global.

\subsection{Mapeo Global: mdx-components.jsx}

Archivo requerido por Next.js para configurar el mapeo global de componentes en archivos MDX:

\begin{lstlisting}[language=JavaScript]
// src/mdx-components.jsx
import { SmartTable, CodeBlock, SmartLink,
         SmartFigure, Solution, Quiz, QuizQuestion,
         MultiColumn, Iframe } from '@/components/mdx';
import { H1, H2, H3, Paragraph, Ul, Ol } from '@/components/mdx/text';

export function useMDXComponents(components) {
  return {
    h1: H1, h2: H2, h3: H3,
    p: Paragraph, ul: Ul, ol: Ol,
    SmartTable, CodeBlock, SmartLink,
    SmartFigure, Solution, Quiz, QuizQuestion,
    MultiColumn, Iframe,
    ...components,
  };
}
\end{lstlisting}

% =====================================================
\section{Resumen de Componentes}
% =====================================================

\begin{longtable}{lll}
\toprule
\textbf{Componente} & \textbf{Categoría} & \textbf{Propósito} \\
\midrule
\endhead
Navbar & Layout & Navegación principal + tema \\
Breadcrumb & Layout & Migas de pan dinámicas \\
Footer & Layout & Pie de página \\
PageTransition & Layout & Transiciones entre páginas \\
SmartTable & MDX & Tabla con KaTeX \\
CodeBlock & MDX & Código con resaltado \\
SmartLink & MDX & Link inteligente int/ext \\
SmartFigure & MDX & Imagen con lightbox \\
Solution & MDX & Respuesta colapsable \\
Quiz/QuizQuestion & MDX & Quiz interactivo \\
MultiColumn & MDX & Layout de columnas \\
Iframe & MDX & Embed seguro \\
Terminal & MDX & Simulación de terminal \\
SubjectsGrid & Feature/INACAP & Grid de asignaturas \\
SemesterTabs & Feature/INACAP & Tabs de semestres \\
Calendar & Feature/INACAP & Calendario académico \\
CftTable & Feature/Bootcamp & Tabla de módulos \\
ProblemTable & Feature/LeetCode & Tabla de problemas \\
ProblemStats & Feature/LeetCode & Estadísticas LeetCode \\
CourseTable & Feature/E-Learning & Tabla de cursos \\
Card & UI & Tarjeta base \\
CardCarousel & UI & Carrusel de tarjetas \\
TableOfContents & Common & TOC automático \\
ScrollReveal & Common & Animación scroll \\
PaginationControls & Common/Table & Paginación \\
SortableHeader & Common/Table & Encabezado sorteable \\
FilterControls & Common/Table & Panel de filtros \\
CustomMDXProvider & Provider & Contexto MDX + tema \\
\bottomrule
\caption{Catálogo completo de componentes}
\end{longtable}

\vspace{1cm}
\begin{center}
\rule{8cm}{0.4pt}\\[0.3cm]
{\small\color{codegray} SaukCode.cl --- Documentación Técnica Vol. II\\
Catálogo de Componentes}
\end{center}

\end{document}
